% !TEX root = ../algebra_02.tex

\section{Возможные порядки конечного поля. Существование поля из $p^n$ элементов}

Пусть $F$ --- конечное поле. Тогда его характеристика отлична от нуля: $\operatorname{char} F = p > 0$, где $p$ --- простое число.
Поле $F$ можно рассматривать как расширение простого подполя $\mathbb{F}_p = \mathbb{Z}/p\mathbb{Z}$.
Пусть степень этого расширения равна $[F : \mathbb{F}_p] = n$. Тогда как линейное пространство над $\mathbb{F}_p$ поле $F$ изоморфно $\mathbb{F}_p^n$. Следовательно, количество элементов в конечном поле равно $|F| = p^n$. Возможные порядки конечного поля исчерпываются степенями простых чисел.

\begin{Theorem}{Существование поля из $p^n$ элементов}{}
	Для любого простого $p$ и натурального $n$ существует конечное поле порядка $p^n$.
\end{Theorem}

\begin{proof}
	Для построения поля из $p^n$ элементов построим поле разложения многочлена. Пусть $K$ --- поле разложения многочлена $X^{p^n} - X$ над полем $\mathbb{F}_p$. Определим множество $F$ как множество всех корней этого многочлена в $K$:
	\[ F := \{ x \in K \mid x^{p^n} - x = 0 \} \]

	Вычислим формальную производную многочлена $f = X^{p^n} - X$:
	\[ f' = p^n X^{p^n - 1} - 1 = -1 \]
	Так как производная не имеет общих корней с $f$, многочлен $f$ не имеет кратных корней. Следовательно, количество корней в точности равно его степени: $|F| = p^n$.

	Покажем, что $F$ является полем.
	Пусть $x, y \in F$, а $\Phi(x) = x^p$ — эндоморфизм Фробениуса:
	\begin{enumerate}
		\item $(x+y)^{p^n} = \Phi^n(x+y) = \Phi^n(x) + \Phi^n(y) = x^{p^n} + y^{p^n} = x + y$, откуда $x+y \in F$.
		\item $(xy)^{p^n} = x^{p^n}y^{p^n} = xy$, откуда $xy \in F$.
		\item Очевидно, что $1 \in F$.
	\end{enumerate}
	Таким образом, $F$ --- подполе в $K$, состоящее из $p^n$ элементов. Так как $K$ порождается корнями $f$, то $K = \mathbb{F}_p(F) = F$. Существование поля из $p^n$ элементов доказано.
\end{proof}

\begin{Proposition}{О примитивном элементе}{}
	Пусть $|F| = p^n$. Тогда существует элемент $a \in F$ такой, что $F = \mathbb{F}_p(a)$.
\end{Proposition}

\begin{proof}
	Мультипликативная группа конечного поля $F^*$ конечна ($|F^*| = p^n - 1$). Известно, что любая конечная подгруппа мультипликативной группы поля является циклической. Следовательно, $F^* = \langle a \rangle$ для некоторого элемента $a \in F^*$. Тогда подполе $\mathbb{F}_p(a)$ содержит все степени $a$, то есть содержит всё множество $F^*$, а значит, $\mathbb{F}_p(a) = F$.
\end{proof}

\begin{Corollary}{Существование неприводимого многочлена любой степени}{}
	Для любого простого $p$ и любого $n \in \mathbb{N}$ существует неприводимый многочлен $f \in \mathbb{F}_p[x]$ степени $n$.
\end{Corollary}

\begin{proof}
	По доказанному выше, существует поле $\mathbb{F}_{p^n}$, и его можно представить как $\mathbb{F}_p(a)$ для некоторого примитивного элемента $a$. Тогда минимальный многочлен элемента $a$ над $\mathbb{F}_p$ является неприводимым и имеет степень $n$.
\end{proof}
